% Front matter for the self-contained Version 3 manuscript.
% The compile-time switch \ifErdosProofClosed is defined by the master file.
% It must remain false until the proof-closure checklist in Appendix A is green
% on one integrated commit.

\begin{abstract}
\ifErdosProofClosed
Let $G_n\sim G(n,1/2)$. We prove that the difference between the chromatic
number $\chi(G_n)$ and the cochromatic number $\zeta(G_n)$ is bounded below by
a positive constant times $n/(\log n)^3$ with probability tending to one. The
argument is uniform across the full rounding phase of the independence-number
threshold. Its first-moment component compares ordinary colorings with signed
cocoloring witnesses supported on four consecutive class sizes. For the second
moment, we derive an exact sign-compatibility identity, encode the large
overlap cells by a matching with local deficits, sum the corresponding physical
matching fibers exactly, and control the residual even-subgraph factor by
injective restriction outside the exposed matching. A bounded-differences
argument then amplifies the resulting positive-probability seed to a
high-probability cocoloring.
\else
This verification draft assembles a self-contained proof architecture for a
full-sequence lower bound on $\chi(G_n)-\zeta(G_n)$ in $G(n,1/2)$. It includes
the complete manuscript argument, the replacement high-skeleton and residual
attachment sections, and an explicit map to the current Lean formalization.
The displayed target theorem is not promoted as proved until the remaining
phase, partial-diagonal, and global assembly gates listed in Appendix~A are
validated on one integrated commit. This status distinction is part of the
manuscript rather than being left only in repository metadata.
\fi
\end{abstract}

\maketitle

\ifErdosProofClosed\else
\begin{center}
\fbox{\parbox{0.91\linewidth}{\small
\textbf{Verification draft.}
The finite Section~8 bridge through the realized endpoint-table deficit sum is
kernel-checked privately, but the full theorem is still conditional on the
remaining formal gates stated in Appendix~A. The publication switch in the
master file must remain disabled until those gates close.}}
\end{center}
\fi

\section*{Introduction}
\addcontentsline{toc}{section}{Introduction}
\label{sec:introduction-v3}

All graphs in this paper are finite, simple, and undirected. A
\emph{cocoloring} of a graph $G$ is a partition of $V(G)$ into nonempty
classes, each of which is either an independent set or a clique. The least
number of classes in such a partition is the \emph{cochromatic number}
$\zeta(G)$. We write $\chi(G)$ for the chromatic number. These are the
standard graph invariants; the signed objects introduced below are auxiliary
counting witnesses and do not define a new invariant.

Let $G_n\sim G(n,1/2)$ be the labeled random graph on $[n]$ in which the
$\binom n2$ possible edges occur independently with probability $1/2$. An
event holds \emph{with high probability} if its probability tends to one as
$n\to\infty$.

Erd\H{o}s and Gimbel asked whether
\[
  \chi(G_n)-\zeta(G_n)\longrightarrow\infty
\]
with high probability \citep[p.~263]{erdos-gimbel-1993}. The problem was
later restated by Gimbel \citep[Section~7.4]{gimbel-2016} and is cataloged as
Erd\H{o}s Problem~625 \citep{bloom-erdos625}. The intended final result is the
following full-sequence quantitative statement.

\ifErdosProofClosed
\begin{theorem}\label{thm:main-v3}
For $G_n\sim G(n,1/2)$,
\[
  \mathbb P\!\left(
    \chi(G_n)-\zeta(G_n)
    \ge
    \frac{(\log 2)^2}{8}
    \log\!\left(\frac{1000}{639}\right)
    \frac{n}{(\log n)^3}
  \right)
  \longrightarrow 1.
\]
In particular, $\chi(G_n)-\zeta(G_n)$ tends to infinity with high probability
along the full sequence of integers $n$.
\end{theorem}
\else
\begin{theorem}[Target theorem]\label{thm:main-v3}
Assume the concrete phase package, the signed four-size first-moment assembly,
the complete partial-diagonal estimate, and the global attachment assembly
listed in Appendix~A. Then, for $G_n\sim G(n,1/2)$,
\[
  \mathbb P\!\left(
    \chi(G_n)-\zeta(G_n)
    \ge
    \frac{(\log 2)^2}{8}
    \log\!\left(\frac{1000}{639}\right)
    \frac{n}{(\log n)^3}
  \right)
  \longrightarrow 1.
\]
\end{theorem}
\fi

The full-sequence quantifier is essential. Near a jump of the
independence-number threshold, changing the relevant class size by one changes
the feasible profile. A result proved only on a density-one set of phase
values therefore need not extend to every integer $n$. Every asymptotic
estimate below is stated uniformly in the complete phase parameter, including
sequences approaching either endpoint of a phase interval.

\subsection*{Background}

The chromatic number of dense random graphs has been studied since the work of
Grimmett and McDiarmid \citep{grimmett-mcdiarmid-1975}. The first-order
asymptotic was established by Bollob\'as \citep{bollobas-1988}; later
refinements include \citet{mcdiarmid-1990},
\citet{panagiotou-steger-2009}, and \citet{heckel-2018}. The cochromatic
number belongs to the theory of generalized chromatic numbers for hereditary
graph properties developed by Scheinerman and Bollob\'as--Thomason
\citep{scheinerman-1992,bollobas-thomason-1995}.

For the chromatic--cochromatic difference, Heckel and, independently, Steiner
obtained the first quantitative evidence toward divergence
\citep{heckel-2024-question,steiner-2024}. Heckel subsequently proved a much
larger lower bound on a phase-dependent set containing approximately $95\%$ of
the integers \citep{heckel-2025-difference}. The remaining difficulty is to
control the complete phase uniformly. The present argument uses the signed
first-moment gain and rare-seed amplification from that work, but replaces the
pairwise sign bound by an exact overlap identity and treats every phase with
one four-size profile.

\subsection*{Contributions}

The proof is organized around five results that are mathematically distinct
and should be read separately.

\begin{enumerate}
\item A phase-uniform comparison between the ordinary coloring root and a
four-size signed cocoloring root gives a separation of order
$n/(\log n)^3$.
\item The sign compatibility of two witnesses is evaluated exactly. The
result is a product of local cell rewards multiplied by the cardinality of a
binary cycle space.
\item A completion-free high-skeleton theorem identifies every
matching-supported physical demand fiber with a product of literal one-cell
partial matching fibers. The ambient falling-factorial loss is then charged
once, after the fiber sum.
\item An exact regrouping identity transports full-support reference weights
through the realized endpoint-table decomposition. The main proof uses a
coarse common charge, while a sharper table-preserving version is retained as
a separate reusable statement.
\item Deletion of the exposed matching is injective on the family of even
residual edge sets. This gives a direct restriction-product estimate and
removes the need for a simple-cycle or walk-kernel decomposition in the final
argument.
\end{enumerate}

The exact signed-overlap identity, the completion-free matching-fiber
factorization, the weighted realized-table regrouping, the restriction-product
lemma, and the rare-seed amplifier are stated as independent propositions or
lemmas. This separates the reusable finite mathematics from the phase-specific
random-graph estimates.

\subsection*{Proof structure}

The proof has three layers.

The \emph{location layer} compares the ordinary first-moment root with the
signed four-size root. The \emph{existence layer} proves a normalized
second-moment estimate for one tangent-rounded signed profile. The
\emph{amplification layer} converts the positive-probability signed witness into
a high-probability upper bound for $\zeta(G_n)$ and intersects it with the
chromatic lower bound.

The difficult overlap calculation is itself divided into four operations:
partial diagonals, canonical high cells, residual attachments, and final
normalization. Each operation has a separate theorem contract. In particular,
no residual local or cycle-space factor is charged in the high-skeleton
section, and no high-cell multiplicity is charged again in the residual
section.

\subsection*{Organization}

Section~1 fixes the phase notation and elementary inequalities. Sections~2--4
locate the ordinary coloring root and derive an unrestricted lower bound for
$\chi(G_n)$. Section~5 constructs the signed four-size profile and proves the
root separation. Section~6 gives the exact signed-overlap identity, and
Section~7 controls partial diagonals. Section~8 proves the completion-free
high-skeleton estimate and endpoint transport. Section~9 treats the residual
attachment by matching restriction. Section~10 proves the rare-seed amplifier,
and Section~11 assembles the final events and the strengthened constant.
Appendix~A records the exact paper-to-Lean map and the remaining verification
gates.
